\section{The induced prior on function neighborhoods}
For a bounded function $f$ on $\mathcal X$ and $0<\delta<1$, write
\[
 B_\delta(f)=\{w:\|z_w-f\|_{\infty,\mathcal X}\le\delta\}.
\]
These events are measurable because the network map is continuous in parameters in the uniform norm on the compact input domain.

\begin{theorem}[Two-sided finite function-prior bounds]
\label{thm:prior}
For every $N\ge1$ and $\gamma\ge1$,
\begin{equation}
 -\frac1\omega\log\pi_{N,\gamma}(B_\delta(f))
 \ge\frac12\VarPath_{H,N,\delta}(f)-\frac{\log2}{2}.
 \label{eq:prior-lower}
\end{equation}
If $p=\VarPath_{H,N,\delta/2}(f)<\infty$, then
\begin{equation}
 -\frac1\omega\log\pi_{N,\gamma}(B_\delta(f))
 \le2p+1+2\log\frac{8(H+1)(p+1)}\delta+2\log2.
 \label{eq:prior-upper}
\end{equation}
In particular the precision-normalized negative log prior is controlled by minimal approximate path variation, with only the displayed logarithmic accuracy and depth terms. No learner-width or input-dimension factor is suppressed.
\end{theorem}
\begin{proof}
Support conditioning leaves the first layer and readout Gaussian and independent, so
\begin{equation}
 \omega Q_N\sim\chi^2_{N(d+1)},\qquad
 \frac{N(d+1)}\omega\le1.
 \label{eq:chi-square}
\end{equation}
Any network in $B_\delta(f)$ has $P(w)\ge\VarPath_{H,N,\delta}(f)$ and $Q_N\ge2P(w)$ by~\eqref{eq:energy-lower}. Markov's inequality with
$\E e^{\omega Q_N/4}=2^{N(d+1)/2}$ yields
\[
 \pi(B_\delta(f))\le
 \exp\left[-\frac\omega2\VarPath_{H,N,\delta}(f)
                      +\frac{N(d+1)\log2}{2}\right],
\]
which proves~\eqref{eq:prior-lower}. An infinite infimum means the event is empty.

For the upper bound, choose any width-at-most-$N$ network $g$ within $\delta/2$ of $f$, with path variation $P$. Apply Lemma~\ref{lem:comparison} at
$\eta=\delta/[8(H+1)(P+1)]$. Its uniform-error $L^2$ norm is at most $\delta/4$, so its comparison law gives probability at least $3/4$ to $B_{\delta/2}(g)\subset B_\delta(f)$. The binary relative-entropy inequality implies, for any event $A$ with $\rho(A)\ge1/2$,
\[
 \log\frac1{\pi(A)}\le2\KL(\rho\|\pi)+2\log2.
\]
Use~\eqref{eq:comparison-entropy}, divide by $\omega\ge1$, and let $P$ decrease to the infimum $p$. This proves~\eqref{eq:prior-upper}.
\end{proof}

The different radii in the two bounds are intentional: approximation at half the requested accuracy leaves room for an open neighborhood of the comparison implementation. The theorem is about function neighborhoods, not positive mass on one exact continuous-domain function. It gives bounds on prior preference; it does not assert an exact monotone ordering between every pair of functions.
